\section{Data}
\label{sec:data}

\paragraph{Venue and ground truth.}
We draw submissions from a recent edition of a major machine learning conference
hosted on OpenReview (primary cohort: \mbox{ICLR}~2026). The venue is chosen for
two properties. First, it operates open review, so \emph{rejected}
submissions, their reviews, and the final decision are all public; venues that
publish only accepted papers cannot support the central claim of this study,
which concerns the reject class. Second, its decisions and reviews were released
in January~2026, \emph{after} the grading model's August~2025 knowledge cutoff,
so the model cannot have memorized the outcome (\S\ref{sec:discussion}); the
earlier, fully pre-cutoff \mbox{ICLR}~2025 cohort is retained only as a
contaminated contrast. For each submission we record two ground-truth signals:
\begin{itemize}\itemsep1pt
  \item the \textbf{decision tier} on the ordered ladder
  reject $<$ poster $<$ oral, and its binary collapse into
  accept/reject; and
  \item the \textbf{mean reviewer rating}, a continuous signal that is richer
  than the discrete decision and partially averages out individual reviewer
  noise.
\end{itemize}
We grade the \textbf{submitted} version of each manuscript---the PDF a reviewer
actually saw---never the camera-ready, so the score is computed on the same
artifact the human signal responded to.

\paragraph{Sampling.}
The natural acceptance rate would starve the accepted tiers, so we sample
\emph{balanced} across decision tiers and report metrics both as sampled and
re-weighted to the venue's natural prevalence (\S\ref{sec:results},
Appendix~\ref{app:robust}). Table~\ref{tab:sample} gives the realized counts.
Sampling is reproducible from a committed cohort manifest, frozen before any
frontier-model grading so the expensive budget is spent on exactly the
pre-registered submissions.

\paragraph{Exclusions (pre-registered).}
We exclude desk-rejected and withdrawn-before-review submissions (no completed
reviewer signal), submissions whose PDF failed to parse, and submissions the
pipeline's self-identity step flags as an already-published manuscript under a
different title (an arXiv-twin leakage risk for the citation audit). Excluded
counts and reasons are reported and the excluded set never enters a primary
metric.

\paragraph{Nested cohorts.}
To establish the relationship at scale while respecting a fixed frontier-model
budget, gradings are organized into two strictly nested cohorts
(Table~\ref{tab:sample}): every sampled submission is graded by the full pipeline
on a cheap model (cohort $M$, full-mini); a stratified subset is additionally
graded by the same full pipeline on a frontier model (cohort $H$, full), with
$H \subseteq M$. The two configurations differ only in the reviewer model, so each
frontier grading has a paired cheap-model score, enabling the mini-to-frontier
bridge analysis of \S\ref{sec:results}.

\begin{table}[t]
  \centering
  \caption{Realized sample. Cohort $M$ (full-mini, cheap model) carries the
  statistical power; the nested frontier cohort $H$ (full) carries the production
  model. Both run the full two-pass pipeline. Counts exclude pre-registered
  exclusions.}
  \label{tab:sample}
  \input{tables/tab_sample}
\end{table}

\paragraph{Data contract and release.}
All analysis reads a two-table contract (one row per submission for labels, one
row per grading run for scores; full schema in Appendix~\ref{app:schema}). We
release the de-identified scored table and the analysis code so every figure and
number reproduces with a single command. We do not publicly release the full
AIPR review text for every identifiable submission by default; that corpus is a
controlled audit artifact that can be provided to editors or reviewers if needed.

\paragraph{PDF provenance and the cohort manifest.}
For every retained submission we record the SHA-256 of the graded PDF
(\texttt{pdf\_sha}) and grade the OpenReview \emph{submitted} manuscript---the
double-blind version reviewers saw at decision time, with no author block
(\texttt{is\_anonymized} is uniformly~1 at this venue)---not the camera-ready or
post-rebuttal revision, pinning the artifact to the version available at the
decision. The realized cohort manifest---submission ids, sampling stratum, and
assigned configuration---is the released \texttt{submissions} and \texttt{gradings}
tables themselves, anchored to the public pre-registration tag; the analysis plan
was frozen before any score was joined to an outcome, so the release \emph{is} the
manifest a reader audits against rather than a separately asserted artifact.
